The starting puzzle is that silence about energy performance commands a price premium relative to the worst certified scores in both property types. Figure~\ref{fig:price_levels} plots the volume-weighted average of market-level hedonic price indices per square metre for each EPC disclosure category from January 2019 to June 2024, where each index holds property characteristics constant within a local market so that comparisons across disclosure categories reflect differences in the disclosure decision rather than differences in property type or location. For apartments (panel~(a)), non-disclosing listings occupy the highest price tier across the \emph{entire} sample window: the ordering Non-disclosed $>$ Efficient (A/B/C) $>$ Energy-intensive (E/F/G) is visible as early as 2019, well before the reform, and the gap widens sharply after July 2021. For houses (panel~(b)), the pre-reform picture differs: non-disclosers begin near the bottom of the price distribution and rise toward the middle post-reform, reflecting the higher baseline audit coverage in the house segment documented in Section~\ref{ss:fact3}. The pre-reform level premium for apartments is consistent with the concentration of non-audited stock in high-price urban markets. Figure~\ref{fig:fact1_premium} isolates the \emph{change} in the gap: it plots difference-in-differences estimates of the hedonic-controlled log price gap between non-disclosing and energy-intensive listings, purging observable composition differences and measured relative to the pre-reform baseline. Because the estimates are normalised to that baseline, the pre-reform coefficients are informative about trends rather than levels. They average $+0.002$ log points for apartments and $+0.008$ for houses and display no discernible drift, indicating that the two series evolved in parallel before July 2021. This is precisely the identifying assumption the design requires; it is not a claim that non-disclosing and energy-intensive listings were priced alike, and Figure~\ref{fig:price_levels} shows that for apartments they were not. After July 2021, when DPE certificates became legally opposable under \emph{loi Climat et R\'esilience}, the estimates shift abruptly upward and stabilize around $+0.044$ log points for apartments and $+0.037$ log points for houses by mid-2024. The reform therefore does not create the level difference so much as widen it sharply, in both the raw series and the composition-adjusted estimates. Classical unraveling predicts the opposite sign: silence should signal bad news and trade \emph{below} verified low-quality listings.
